\chapter{Sentinel Lymph-Node Biopsy}

\section*{Introduction \& Clinical Indications}
Sentinel lymph-node biopsy identifies and removes the first draining lymph node or nodes from a primary tumor to stage regional disease. It is commonly used in selected breast cancer and melanoma patients and in other cancers according to disease-specific guidance. A negative sentinel node may allow avoidance of a more extensive nodal dissection, reducing morbidity, but indications depend on tumor type, stage, imaging, biology, prior treatment, and whether the result will change management.

\section*{Relevant Surgical Anatomy}
Lymphatic drainage varies by skin site and breast location. The axilla, internal mammary chain, groin, and neck contain important vessels, nerves, and lymphatic channels. Mapping may identify one or more sentinel nodes or an unusual drainage basin. The intercostobrachial, long thoracic, thoracodorsal, femoral, and cutaneous nerves may be at risk depending on the basin.

\section*{Preoperative Workup \& Preparation}
Pathology, imaging, tumor size, clinical nodal examination, and prior surgery or radiation are reviewed. Mapping uses radiotracer, blue dye, magnetic tracer, or combinations. Allergy history, pregnancy, anticoagulants, and anesthesia risk are assessed. Consent covers failed mapping, false-negative results, bleeding, infection, nerve injury, allergic reaction, lymphedema, seroma, and possible additional surgery.

\section*{Surgical Technique \& Procedure}
The tracer is injected according to the tumor and protocol. A probe, visual dye, or magnetic detector identifies the draining node. The node or nodes are removed through a targeted incision and sent for permanent pathology, with intraoperative assessment in selected settings. The wound is closed after hemostasis; completion axillary or nodal dissection is not automatic.

\section*{Complications \& Risks}
Risks include pain, bruising, seroma, infection, allergic reaction to dye, temporary or persistent numbness, lymphedema, shoulder or limb stiffness, and false-negative or unsuccessful mapping. Pathology may require further surgery or systemic or radiation treatment.

\section*{Postoperative Care \& Recovery}
Most patients resume light activity quickly. Wound care, pain control, arm or limb movement, and lymphedema education are provided. Fever, expanding swelling, drainage, severe pain, or new limb swelling requires review. The final pathology report guides the multidisciplinary treatment plan.

\section*{Outcomes \& Prognosis}
The procedure provides staging information rather than treating every microscopic cancer cell. Prognosis depends on the primary tumor, sentinel-node pathology, molecular features, and adjuvant therapy. A negative node does not eliminate distant recurrence risk, and a positive node does not automatically require complete nodal dissection in every modern treatment pathway.

\section*{References}
\begin{enumerate}
  \item American Society of Clinical Oncology. Sentinel Lymph-Node Biopsy Guideline.
  \item American Society of Breast Surgeons. Consensus Guideline on Sentinel Lymph-Node Biopsy.
  \item National Comprehensive Cancer Network. Breast Cancer and Melanoma Guidelines.
\end{enumerate}

\clearpage
